\documentclass[11pt]{article}
\PassOptionsToPackage{numbers}{natbib}
\usepackage[review]{acl}
\makeatletter\acl@anonymizefalse\makeatother
\usepackage{times}
\usepackage{latexsym}
\usepackage[T1]{fontenc}
\usepackage[utf8]{inputenc}
\usepackage{microtype}
\usepackage[htt]{hyphenat}
\usepackage{booktabs}
\usepackage{amsmath}
\usepackage{graphicx}
\sloppy
\setlength{\tabcolsep}{4pt}
\title{Team Name at SemEval-2027 Task 1: Multi-evidence re-ranking with a frozen first stage}
\author{Author Name \\ Affiliation \\ \texttt{email@example.org}}
\date{}
\begin{document}
\maketitle
\begin{abstract}
RETECO Track 1a asks for documents that answer a \emph{complete} question about a historical period: long queries whose temporal and exclusion conditions sit in the body rather than the title, and judgements that mark \textbf{several} documents per question (3.44 on average, over \textbf{2.34 annotated steps}). We describe a two-stage ranker - a zero-shot dense retriever plus a windowed cross-encoder fine-tuned on the 1,071 training queries that carry at least one positive among their dense top-100 candidates - that scores 0.3581 domain-macro nDCG@10 on the 13-domain dev set, against 0.2943 for the dense stage and 0.1147 for the organisers’ BM25 baseline. Against the same pipeline with the cross-encoder left at zero-shot (0.3249), the fine-tuned reranker contributes +0.0332 (domain-clustered 95\% CI [+0.0192, +0.0472]) at a fusion form fixed before fine-tuning, and reproduces at +0.0326 / +0.0279 under two further seeds. The remaining headroom is only partly recall: at the 100 candidates the reranker sees, 80.7\% of annotated steps have at least one relevant document (96.3\% within the dense top-1000), yet the fused top ten covers only 59.5\% of the annotated steps and the cross-encoder’s own top ten 51.4\%. Ten temporal probes, one strictly pre-registered test (+0.0039 against a +0.02 bar) and a fixed year-consistency reordering that \emph{lowers} nDCG@10 in 7 of 8 domains do not close it. Supervising every annotated step to have a positive ranked high is a wash (0.3593 and 0.3600 against 0.3581, both intervals include zero) and leaves coverage where it was. Restoring the 140 training queries with an empty candidate pool is a wash too (0.3565). Only the constraint-faithful 64-word query pays a price: 0.3443, 0.0138 below the frozen system. Every intervention in the six-step ladder is neutral or negative on dev, and the frozen system is kept. A bounded replacement stage is the exception: letting a 4B judge compare three external candidates against the documents at ranks 8-10, twice with the order swapped and at most one swap per query, raises the same system to 0.3646 (+0.0065, 95\% CI [+0.0021, +0.0117]) against a ceiling of 0.3769 for that candidate space, while adding document context to the same prompt does nothing and an off-the-shelf reranker given the identical pool loses.

\end{abstract}
\section*{1 Introduction}
Track 1a gives 13 independent corpora, long queries, binary qrels and no timestamps. Our entry has two parts: a strong, reproducible zero-shot-to-supervised pipeline, and a diagnosis of why the obvious temporal fixes do not work.

Contributions: (i) a baseline ladder over all 13 domains under the official domain-macro statistic, with paired bootstraps over queries and over domains; (ii) a supervised cross-encoder that adds +0.0332 without tuning the fusion weight on dev; (iii) a multi-evidence analysis showing that the system’s failure mode is "one easy document" inside a candidate pool that is itself incomplete; (iv) negative evidence on explicit temporal modelling: one pre-registered test and a fixed year-consistency rule that fails; (v) a bounded, pre-registered replacement protocol with its exact oracle ceiling and a costed comparison against per-document scoring, an off-the-shelf reranker and a PRP-style policy; (vi) a decomposition of why the cross-encoder with the best standalone score is not the one that fuses best - anchoring, attenuation and score calibration - and what that implies for using a single-stage metric to select checkpoints.

\emph{Relation to prior work.} The reranker is the standard retrieve-then-rerank cross-encoder \citep{nogueira2019passage}, initialised from M3-Embedding \citep{chen2024m3}. Our multi-evidence objective is a sibling of listwise and LambdaRank training \citep{burges2010lambdarank,wang2018lambdaloss}, with one difference: instead of ordering one list, it asks that every official step of a question own a document in the top ten. Multi-stage fusion has its own line - list-aware reranking of the first stage’s list \citep{zhang2022hlatr} and joint selection of documents that cover the answer set \citep{mao2021jpr} - but both learn the second stage jointly with the list it consumes (HLATR adds a list-aware module over retrieved and reranked lists; JPR reranks the set autoregressively), whereas we freeze the first stage and ask what one second stage can add at a fixed fusion weight. The contribution is the measured behaviour of that frozen system, not a new fusion architecture. Temporal retrieval is the benchmark’s own subject \citep{abdallah2026tempo}; our contribution there is negative evidence about what does not transfer from the two-domain pilot to the full suite.

\section*{2 Task, metric, data}
\textbf{Metric.} Official aggregation is domain-macro nDCG@10: per-query nDCG@10, averaged within domain, then over the 13 domains, so a 3-query domain weighs as much as a 240-query one, and queries missing from a run score zero. dev = 519 queries. 26.6\% of dev questions receive exactly zero from the frozen system, and 37.7\% of those have no relevant document in the top 100 at all. Those zeros dominate the statistic: had the 138 zero questions merely reached half their domain’s average, domain-macro would be 0.4283 instead of 0.3581 (full average: 0.4985, \texttt{results/ZERO\_HEADROOM.json}).

\textbf{Data.} 1,167,159 documents, 1,211 train queries, binary qrels (so nDCG@10 is discounted precision at 10). Per-domain counts in \texttt{results/corpus\_manifest.json}. Queries carry a title and a body separated by \texttt{<p>}; the body holds the temporal range, exclusions and comparisons (Section 5.3).

\section*{3 System}
\textbf{Dense first stage.} \texttt{mxbai-embed-large-v1}, full query with the model’s retrieval instruction, mean pooling, cosine similarity, fp16, \texttt{top-100} candidates per query. Dev: 0.2943 (0.2956 with CLS pooling: a wash).

\textbf{Windowed cross-encoder.} \texttt{bge-reranker-v2-m3} (568M). Documents are split into 250-token windows with stride 192, each window is scored against the query, and the document score is the maximum over its windows, i.e. a document is as relevant as its best passage; query capped at 256 tokens. Pooling was fixed a priori and is not ablated (Section 7). Zero-shot at fusion 1:1: 0.3249.

\textbf{Fusion.} Per-query z-normalisation of the dense and cross-encoder lists, combined at 1:1. The weight was fixed on the training split before any dev inspection; Section 5.1 shows the choice is not what produces the gain.

\textbf{Supervised fine-tuning.} Listwise softmax over a 20-document group sampled from each training query’s 100 candidates - 21,420 query-document pairs in 1,071 groups, the log-sum-exp running over the group’s positives - 8 of 24 layers frozen, lr 1.5e-5, seed 17, 2 epochs; epoch 1 was kept, by a 0.3029-against- 0.3012 internal-validation margin. Checkpoint selection used the cross-encoder’s own ordering on 45 held-out training queries, not the fused system; that run predates domain-macro logging and selected on the plain mean, and on every later arm the two criteria pick the same epoch. Frozen dev: \textbf{0.3581}.

\textbf{Multi-evidence (step) supervision.} Track 1a gives one relevance judgement per query, but the organisers’ qrels also carry \emph{step} annotations: a question that needs several pieces of evidence is decomposed into steps, and each step lists the documents that satisfy it. We use them only in training; at test time the ranker still sees the full question and one document. With \texttt{P\_j} the step-\texttt{j} positives inside a query’s candidate list, the ordinary listwise term \texttt{-log sum\_\{d in P\} p(d)} is dominated by the single easiest positive, so we add

\texttt{`} L\_step = -(1/m) * sum\_j log sum\_\{d in P\_j\} p(d),    L = L\_list + lambda * L\_step \texttt{`}

which forces \emph{every} step to have a positive ranked high instead of letting one easy document carry the query. Steps whose positives are absent from the candidate list are skipped (a step is covered for 91.2\% of step instances under our candidate construction); \texttt{m} counts the covered steps. \texttt{lambda} is not tuned on dev: two values (0.3, 1.0) are trained on the same groups, and both raise the train-internal split above the baseline (0.2882 and 0.2913 against 0.2865), yet both are rejected on dev (S5.6).

\section*{4 Results}
\begin{table}[t]
\centering
\footnotesize
\resizebox{\columnwidth}{!}{%
\begin{tabular}{lrrr}
\toprule
system & dev domain-macro nDCG@10 & $\Delta$ vs zero-shot z1:1 & domain-clustered 95\% CI \\
\midrule
official BM25 & 0.1147 & --- & --- \\
dense (mxbai, top-100) & 0.2943 & --- & --- \\
dense + zero-shot CE, z 1:1 & 0.3249 & --- & --- \\
+ SFT CE, 1 epoch & 0.3336 & +0.0087 & [+0.0037, +0.0142] \\
\textbf{+ SFT CE, 2 epochs (frozen)} & \textbf{0.3581} & \textbf{+0.0332} & \textbf{[+0.0192, +0.0472]} \\
same recipe, seed 23 & 0.3575 & +0.0326 & [+0.0184, +0.0470] \\
same recipe, seed 29 (two independent reruns) & 0.3528 / 0.3527 & +0.0279 & [+0.0160, +0.0405] \\
\bottomrule
\end{tabular}}%
\end{table}

The gain holds when the three smallest domains are dropped (0.3267 against 0.2985 on domains with n >= 19 queries) and in 11 of 13 domains; the two exceptions, bitcoin and quant, are flat (-0.0023 and -0.0009). Seed spread over three seeds is 0.0029, an order of magnitude below the effect, while a fourth run of the same recipe through the current code path scores 0.3575: the retraining band is 0.3527-0.3581 and the frozen artifact sits at its top (S5.5). Rank-fusion forms were screened too and are reported in S5.1, where the RRF list that out-of-fold selection prefers is worse than the z 1:1 fusion used here.

\section*{5 Analysis}
\subsection*{5.1 Is the gain an artefact of the fusion weight?}
No. A sweep over nine weights (0.5 to 2 on the cross-encoder, z-normalised) is positive everywhere (+0.0089 to +0.0346; the w=1 row reproduces the headline within bootstrap noise), and refitting the weight out-of-fold on 45 held-out training queries selects a dense-heavy RRF 1:4 that is \emph{worse} than z 1:1, so the gain is not bought by tuning on dev.

\subsection*{5.2 The bottleneck is multi-evidence coverage over an incomplete pool}
We measured, per system and without extra annotation, how much of the official step structure reaches the top ten (\texttt{scripts/dev\_multievidence.py}):

\begin{table*}[t]
\centering
\scriptsize
\resizebox{\textwidth}{!}{%
\begin{tabular}{lrrrrr}
\toprule
system & nDCG@10 & relevant docs in top-10 & questions with >=2 & step coverage of top-10 & steps retrievable @1000 \\
\midrule
dense & 0.2943 & 1.03 & 31.8\% & 51.3\% & 96.3\% \\
+ zero-shot CE & 0.3249 & 1.11 & 35.6\% & 54.9\% & 96.3\% \\
+ SFT CE (frozen) & 0.3581 & 1.21 & 38.3\% & 59.5\% & 96.3\% \\
\bottomrule
\end{tabular}}%
\end{table*}

Supervision improves every column but leaves most of the structure unused. At depth 1000 seven of ten steps are retrievable and still four of ten never reach the top ten; at the depth the reranker sees (100 candidates per query) only 80.7\% of annotated steps are reachable and 88.3\% of questions keep a positive, so a fifth of the step structure is lost before the cross-encoder runs. Truncating the fused list from 1000 to 100 rows leaves its top ten unchanged (59.45\%, \texttt{results/MULTIEVIDENCE\_sft\_e2bc\_sub100.json}), so depth moves reachability only; the 51.4\% is the cross-encoder’s own ordering over those 100 candidates. The retrieval floor is small but visible: in 10 of 519 dev questions (1.9\%) no relevant document appears anywhere in the dense top-1000, and 52 of the 138 questions the frozen system scores at zero (38\%) have none in the top 100, where no reranker can lift them. Coverage tracks nDCG@10 (Spearman 0.69; all-covered questions average 0.595 against 0.187, \texttt{results/STEP\_COVERAGE\_CORR.json}), though partly mechanically, so we read it as a check rather than a result.

The gain does not concentrate on the questions that need several evidence documents: on the 151 two-step dev questions it is +0.0200 (CI [-0.0033, +0.0463]) against +0.0711 on single-step ones (Appendix A.4).

\subsection*{5.3 Query truncation drops the conditions}
In 1,730 train+dev queries, 44.8\% (46.4\% of dev: 241 of 519) carry at least one condition (year 14.3\%, temporal relation 19.0\%, exclusion 16.3\%, comparison 14.0\%) that appears \textbf{only in the body}, so it never reaches the cross-encoder, which is trained on the title alone; the dense stage reads the full question. That truncation is not, by itself, what the ranker is missing: on the 241 condition-lost queries the frozen system scores \textbf{0.3656 against 0.3518} on the rest, and its gain over zero-shot CE is about twice as large there (+0.0481 vs +0.0246), although the dense stage already reads the full question and only the reranker does not. A same-budget body-visible composer therefore has to \emph{earn} the claim, so we pre-register that its gain must concentrate on the lost subset. Two 64-word composers were built - a cue-sentence selector (title plus only the body sentences carrying a condition) and the plain full question (title plus the first 64 body words). On the 45-query internal split no representation wins outside noise: title 0.2827, full 0.2869, selector 0.2744, fixed body 0.2703, frequency summary 0.2614 (\texttt{results/QVARIANT\_WINNER.json}); the \textasciitilde{}+-0.03 resolution of that split (S7) makes dev the decisive test, and the selector - the more faithful of the two - was the worst, so dev carries the \emph{generous} arm. Trained end to end on the 64-word full question it scores \textbf{0.3443} against 0.3581 (Delta -0.0138, query-clustered CI [-0.0272, -0.0005], domain-clustered [-0.0338, +0.0052]) and loses on both subsets: 0.3566 vs 0.3656 on the 241 condition-lost queries and 0.3296 vs 0.3518 on the other 278. The loss is thus \emph{smaller} where the arm was designed to help (-0.0090 vs -0.0222), consistent with body conditions not being the missing signal, but no subset gains, so the pre-registered bar is not met (\texttt{results/CONSTRAINT\_SUBSET\_cx\_qv.json}).

\subsection*{5.4 Temporal modelling did not help}
Ten probes - year/anchor features, a semantic time-consistency score, oracle and LLM step decomposition, HyDE, temporal hard negatives, a time-aware constraint rerank, a train-fitted gate, a fixed year-consistency reordering, and a coverage diagnostic. Exactly one was strictly pre-registered: it gained +0.0039 ([-0.0023, +0.0103], domain [+0.0006, +0.0142], P=0.893) against a +0.02 bar. The year rule fails cleanly: ordering the frozen candidates by year consistency lowers nDCG@10 in 7 of 8 anchored domains (macro 0.2980 to 0.2498). That is a fixed rule losing, not an upper bound - the reordered list is not the best possible one, and the rule itself is what a temporal ranker would have to beat. In this corpus a document mentioning another period is often still the right answer.

\subsection*{5.5 The rest of the ladder, in brief}
Three further steps are reported in full in Appendix A. \textbf{Loss forms}: averaging the log-loss over all positives (0.2809), an nDCG-weighted permutation loss (0.2764) and domain-sqrt sampling (0.2912) all sit within the 0.010 tie band of the current objective (0.2865) on the 45-query split and were not promoted to dev, so their dev effect is unmeasured. \textbf{Step supervision}: requiring every annotated step to have a positive ranked high gives 0.3593 and 0.3600 against 0.3581 (domain-clustered CIs [-0.0021, +0.0048] and [-0.0014, +0.0058]) - a wash, and coverage stays at 59.5\% / 59.7\% rather than rising. \textbf{Corrections}: restoring the 140 empty-pool queries is also a wash (0.3565, -0.0017 [-0.0088, +0.0061]); the 64-word full question is the one arm that pays (0.3443, -0.0138 [-0.0272, -0.0005]). With S5.3 and S5.4 the message is consistent: no retrained arm shows a reliable gain over the frozen system, and only the query-rewriting arm loses to it.

\subsection*{5.6 A better cross-encoder is not a better system}
The checkpoint screen scores the cross-encoder alone; the system is the dense+CE fusion, and the two rank the arms differently. \texttt{cx\_qv} is the stronger cross-encoder by 0.0387 (0.3152 against 0.2765) and 0.0132 \emph{worse} fused at the frozen weight.

Three mechanisms explain it (519 dev queries, frozen dense top-1000, per-query z-normalised 1:1 fusion). \emph{(i) Anchoring}: the fused top ten is 73-75\% dense documents and only 54-60\% of it is also in the CE top ten, so a changed ordering must displace dense documents to matter. \emph{(ii) Attenuation}: a CE gain transfers at slope 0.29 - +0.2082 on the 169 queries the CE improves becomes +0.0438 fused, and -0.1661 on the 135 it hurts becomes -0.0665. \emph{(iii) Profile}: on the 215 queries where the CE’s own nDCG@10 does not move (|delta| <= 0.01) the fused score still falls 0.0113, and a transplant isolates the two - \texttt{cx\_qv}’s ordering with \texttt{lx\_possum}’s score profile scores 0.3500 (-0.0075) while the reverse falls to 0.3421 (-0.0154).

Re-weighting the fusion narrows the gap but does not close it. The weight is a free parameter and the optimum is model-specific: on the 45 held-out queries \texttt{lx\_possum} prefers 0.75 and \texttt{cx\_qv} 1.25, and carrying those to dev gives 0.3522 and 0.3509 - a 0.0053 loss for the model the weight was tuned for and a 0.0066 gain for the other, still 0.0066 under the frozen system. In the controlled probe, where the z-profile is computed over the same 100 candidates, the sweep settles at 1.0 for \texttt{lx\_possum} (0.3575) and at 1.5-2.0 for \texttt{cx\_qv} (0.3532, 0.3546): at the frozen 1:1 weight the better standalone model is the worse fused one, and even each model’s own best weight leaves it 0.0029 below the frozen arm’s best. A calibration-free RRF (k = 60) gives 0.3433 and 0.3462: it orders the two correctly and trails the frozen system by 0.012-0.015. On those same queries the model-level inversion reproduces with the candidate list fixed (\texttt{cx\_qv} 0.2893 alone against 0.2462, and 0.3197 against 0.3325 fused); within a model both criteria prefer epoch 1, and that preference is itself noise (epoch 1 on dev: 0.3560 against 0.3575, CI [-0.0047, +0.0020], 423 of 519 tied). A single-stage metric is thus a good ranker of cross-encoders and a poor ranker of the system they feed.

\textbf{A general-purpose reranker does not fill the gap.} Qwen3-Reranker-0.6B - 0.6B parameters, close to our 0.57B cross-encoder - was scored on the same dense top-100 under its official input recipe (budgeted document body, think suffix, P(yes)): 0.2313 on its own, below our fine-tuned 0.2765, and 0.2996 when fused at the same 1:1 weights, 0.058 under the frozen system. Candidates and fusion are identical in both runs, so the difference is the task fine-tune rather than the first stage or the fusion rule, and the two models are of comparable size, so it is not a capacity story either.

\section*{6 Reproducibility}
Every number in the tables is recomputed from released artefacts rather than transcribed: per-query score JSON for each arm, the dev TREC runs, and the official qrels. \texttt{scripts/check\_claims.py} recomputes the headline figures from those files and fails if one drifts, and the pipeline itself is three commands - dense retrieval at depth 1,000, reranking at depth 100, then a z-normalised weighted sum with the weights fixed at 1:1. The cost is one consumer GPU (an RTX 4060 Laptop with 8 GB): encoding the 1.17 M documents is the dominant term, about six hours at the 52 documents/s we measure; reranking 519 queries at depth 100 takes about 80 minutes; a two-epoch training run takes tens of minutes. The three seeds and the two checkpoint epochs per run are released as runs, so the retraining band in S4 can be re-derived without a GPU.

\section*{7 Limitations}
Domain-macro lets 3-query domains weigh as much as 240-query ones, so a single-domain regression can move the headline; the temporal findings are about \emph{textual} time only, not about systems that have timestamps; and the supervised results are single-corpus and single-family (one dense model, one cross-encoder). The window/pooling schedule (250-token windows, stride 192, max over windows, 256-token query cap) was fixed a priori and is not ablated; the dense encoder’s pooling was ablated late and is a wash (CLS 0.2956 against 0.2943 for mean); the 1:1 fusion weight and the two-epoch schedule were picked on the training-internal split, so the frozen number carries that selection. The step-supervision results are reported at one seed; a second seed was not affordable within the GPU budget.

Selection is noisy by construction: the 45-query internal split improves the sum-positive arm by +0.0135 with a domain-clustered interval of [-0.0317, +0.0571], i.e. it cannot resolve differences below about 0.04. Every arm here is therefore screened internally and judged on dev with a paired domain bootstrap (half-width 0.007-0.014); dev gains below +0.010 are treated as unresolved. The screen did not promote every loser (S5.5) - what it cannot do is rank these arms: the five it scored lie within 0.006 of each other, dev puts them between -0.0017 and +0.0019 of the frozen system, and every interval includes zero. Retraining is not free: fresh pools and seeds give 0.3575 / 0.3528 / 0.3527 and the current code path 0.3575, so 0.3581 sits at the top of a 0.3527-0.3581 band rather than being a property of the recipe.

\section*{8 Error analysis}
A quarter of dev questions (138/519, 26.6\%) score exactly zero. The split separates cleanly by candidate-pool recall: \textbf{52 (37.7\%) have no relevant document in the dense top-100} - no top-100 reranker can score them - and 86 (62.3\%) have the gold in the pool and lose it to ordering. Zeros cluster where the corpus is thin (history 46, politics 14, bitcoin 13). Three worked examples, one per mechanism (\texttt{results/DEV\_ERROR\_ANALYSIS.md}):

\begin{table}[t]
\centering
\footnotesize
\resizebox{\columnwidth}{!}{%
\begin{tabular}{lrrrr}
\toprule
failure & query & gold & dense rank & nDCG@10 \\
\midrule
pool miss & travel/201562\_1094 & 2 & 656, not retrieved & 0.0000 \\
ordering & history/70583\_271 & 1 & 24 (CE 14, fused 17) & 0.0000 \\
step coverage & history/59309\_621 & 2 & 1, 10 (CE 1, 72) & 0.6131 \\
\bottomrule
\end{tabular}}%
\end{table}

The pool miss is a PNR booking-limit question whose evidence sits outside the 100 candidates the reranker sees (dense rank 656); a deeper second stage could reach that document, but nothing downstream of the reranker can. The ordering case ("How do we know when Nero was born?") has its gold at dense rank 24, moves to 14 for the cross-encoder and 17 fused, and scores zero. The step case ("How many people watched the moon landings after Apollo 11?") keeps the Apollo 11 figures at rank 1 but pushes the comparison document from dense rank 10 to 72 for the title-only cross-encoder: one of two evidence documents survives, nDCG@10 stays at 0.61, and step coverage fails. That shape is common, not anomalous: only 38.3\% of questions place two or more relevant documents in the top ten (S5.2).


\section*{9 Bounded replacement of the retrieval tail}
The frozen system fixes the first ten positions. This section asks a narrower question than re-ranking: can a bounded set of \emph{external} candidates replace one tail document and pay for itself? The protocol was frozen before any judge output was seen. Candidates are the union of the dense top-20 and the full-query cross-encoder top-20, minus the frozen top ten, ordered by their best rank in either source, capped at three; each candidate is compared with the documents at ranks 8, 9 and 10, giving at most nine pairs per query; at most one swap is executed, and the decision threshold is the existing rule's 0, never tuned on dev. Candidates are generated from retrieval scores only - no qrels enter the pipeline; labels are read only to score the resulting ranking.

The ideal score of \emph{exactly this} space is 0.3769 (+0.0188 over the frozen 0.3581); 68 of 519 queries contain at least one beneficial pair and 15 of those sit at zero. The figure of 0.4098 that an earlier, uncapped pool produced is not this ceiling and is not quoted as one.

\begin{table*}[t]
\centering
\scriptsize
\resizebox{\textwidth}{!}{%
\begin{tabular}{lrrrrrr}
\toprule
arm & nDCG@10 & delta & 95\% CI & help/harm/neutral/no-swap & prompts & seconds \\
\midrule
frozen system & 0.3581 & - & - & - & 0 & - \\
Direct (Qwen3-4B, both orders) & 0.3646 & +0.0065 & [+0.0021, +0.0117] & 32/19/273/195 & 9,342 & 873 \\
Context (+ 9 retained excerpts) & 0.3632 & +0.0051 & [+0.0013, +0.0094] & 27/7/219/266 & 9,342 & 1,519 \\
per-document scoring, same model & 0.3598 & +0.0017 & [-0.0046, +0.0078] & 35/18/412/54 & 6,747 & 474 \\
cheap-gap pair selection & 0.3626 & +0.0045 & [+0.0004, +0.0093] & 19/11/168/321 & 2,630 & 249 \\
Qwen3-Reranker-4B, one swap & 0.3616 & +0.0035 & [-0.0032, +0.0105] & 27/20/421/51 & 6,747 & 587 \\
Qwen3-Reranker-4B, rerank the 13 & 0.3173 & -0.0408 & [-0.0628, -0.0188] & 134/216/169 & 6,747 & 571 \\
Qwen3-Reranker-4B, rerank dense-100\textsuperscript{\dag} & 0.2213 & -0.1368 & [-0.1649, -0.1085] & - & 51,900 & 5,523 \\
\bottomrule
\end{tabular}}
\caption{Bounded replacement on all 519 dev queries (paired bootstrap, 10k draws, queries resampled within domain). Prompts count judge calls, not documents. \textsuperscript{\dag}This row was measured with the right-padded implementation described in Section~\ref{sec:verif} and was not re-run; it is reported only to show the direction.}
\end{table*}

\textbf{Only pairwise judging pays.} Direct clears the pre-registered bar of +0.005 with a confidence interval above zero; Context clears it by less than a thousandth and its paired difference against Direct is -0.0015 (95\% CI [-0.0051, +0.0016], p = 0.81) while consuming 77\% more tokens and 74\% more wall time. Direct improves nine of the thirteen domains, leaves three flat and worsens one (bitcoin, -0.0085). Retained-document context therefore buys nothing here, and the paper does not claim it. Scoring the same documents one at a time reaches +0.0017 with an interval that includes zero, and an off-the-shelf 4B reranker given the identical pool reaches +0.0035 with an interval that includes zero: with an official input recipe and the same 112-word query-matched passage, the reranker neither beats the judge nor is broken. Asked to reorder the same thirteen documents it loses four points of nDCG. These are the corrected numbers after the padding and prompt-template fixes described in Section~\ref{sec:verif}; an earlier implementation had made the reranker look far worse, and the earlier conclusion that it is clearly inferior is withdrawn. A cheap-gap pair-selection policy -- order the nine pairs by a dense-plus-cross-encoder gap and judge at most three of them -- reaches +0.0045 with 72\% fewer judge calls than Direct: most of the gain at a quarter of the cost, although the paired Direct-minus-policy interval includes zero, so we present it as a cost-quality trade-off rather than a proven loss.

\textbf{Budget.} Using 0, 1, 2 and 3 candidates (0, 3,114, 6,228 and 9,342 prompts) gives 0.3581, 0.3610, 0.3655 and 0.3646. The second candidate carries the gain and the third is slightly negative on dev; because that ordering was observed after the protocol was frozen, cap 3 remains the reported configuration and cap 2 is reported as a dev-only budget point.
Replaying the same stored predictions while judging only $k$ pairs shows how much of the gain a smaller budget keeps: with a single pair the deltas are $-0.0007$ (frozen order), $+0.0043$ (cheap-gap order), $+0.0037$ (random order) and $+0.0075$ (oracle order); at three pairs they are $+0.0002$, $+0.0046$, $+0.0060$ and $+0.0115$; at all nine every ordering returns +0.0065. Judging \emph{some} pair early is enough to capture most of the effect, and the gap-based order does not beat a random pair at these budgets, so the ordering saves calls rather than adding quality.

\begin{table}[t]
\centering
\scriptsize
\resizebox{\columnwidth}{!}{%
\begin{tabular}{lrrr}
\toprule
decision rule & nDCG@10 & delta & 95\% CI \\
\midrule
bidirectional + abstain (frozen) & 0.3646 & +0.0065 & [+0.0020, +0.0118] \\
no abstain option & 0.3647 & +0.0066 & [+0.0012, +0.0128] \\
forward order only & 0.3596 & +0.0015 & [-0.0071, +0.0096] \\
forward only, no abstain & 0.3623 & +0.0042 & [-0.0009, +0.0095] \\
\bottomrule
\end{tabular}}
\caption{Decision-rule ablations, recomputed from the stored per-order probabilities of the Direct run; the first row reproduces the arm exactly.}
\end{table}

\paragraph{Where the recoverable gain sits.} Of the 4,524 non-relevant slots in the 519 frozen dev top-tens, 2,497 have a relevant document inside the dense top-100 and outside the top ten, and every query can be improved by at least one replacement (mean best single fix $+0.410$ per query). Splitting the recoverable gain by position, 43.6\% of it sits at ranks 1-3, 35.0\% at 4-7 and only 21.4\% in the 8-10 window the frozen protocol can touch, which is why the bounded stage recovers a third of this ceiling and why widening the action space is the next step. A model-assisted reading of the 100-row worksheet (a 4B model labels the condition the incumbent fails; this is model annotation, not human adjudication) attributes 41 of the fixable slots to temporal conditions, 39 to entity scope, 10 to numeric bounds, 4 to exclusions and 3 to comparisons, with 2 where no violation was found - the same temporal-dominance the body-only condition analysis showed, now at the level of individual ranking decisions.

\paragraph{Query reasoning.} The same 4B judge wrote an explicit analysis for every dev query (entities, time or version conditions, exclusions, evidence type, then search keywords). Encoding four query views with the frozen dense encoder gives domain-macro nDCG@10 of 0.2933 (original), 0.2930 (analysis appended), 0.2546 (analysis alone) and 0.2915 (short query alone): the transformation is neutral, and the analysis alone loses four points. The published query-reasoning lift relies on GPT-4o analyses with ReasonIR embeddings; a 4B analyst paired with this retriever does not reproduce it.

\paragraph{Second starting system.} Starting from the dense ranking
(0.2943 nDCG@10, no cross-encoder in the pool) instead of the fused one,
the same candidates, judge and one-swap rule reach 0.3049
(+0.0106, 95\% CI [+0.0048, +0.0167]; 61 helpful, 15 harmful swaps).
The mechanism is therefore not specific to the fused start, and it gains
more where the top ten contains more clearly wrong slots.
Starting from the Qwen3-Reranker-4B ordering of the dense top-20
(0.2565) instead, the same rule reaches 0.2696 (+0.0131, 95\% CI [+0.0075, +0.0190]),
so the gain is not an artefact of the fused start's top ten being right.



\paragraph{Constraint-violation filter.} Asking instead for an absolute verdict -- a yes/no question about whether a document satisfies every explicit condition in the query -- collapses onto relevance: candidates outside the frozen top-ten are rated satisfying by roughly zero percent, so a relative rule fires on 3 of 519 queries, all neutral, and the score is unchanged at 0.3581. The A/B framing is what makes the comparison useful, because it forces a comparison with the incumbent instead of an absolute verdict.

\paragraph{Widening the action space.} Widening the action space does not pay at this judge quality. With candidates from the dense and CE top-100 and any of the ten slots available, the 519 dev queries fall to 0.35725 / 0.35654 / 0.35274 / 0.35239 at budgets of 2 / 6 / 12 / 18 judged actions (deltas $-$0.0009 / $-$0.0016 / $-$0.0054 / $-$0.0057). The degradation grows with the budget because each extra action is another chance for the judge to accept a wrong swap, and a wrong acceptance at rank 1 costs an order of magnitude more than at rank 10 (mean change $-$0.083 at slot 1 versus $-$0.002 at slot 10 over 9,342 judgements) while the judge's reliability does not improve with position. A position-aware expected-gain gate fitted on the leak-free folds (change $\sim$ $1/\mathrm{slot}$, margin, margin$/\mathrm{slot}$; 1,721 fold actions, 154 helpful) accepts 19 swaps on dev and nets 0.000000: it removes the damage but finds no gain. The precision of the small CE-sourced pool, not the breadth of the action space, is what carries this stage.

\paragraph{Allocation versus judgement.} The budget curve shows that widening the space fails, but not why. Judging which action to spend a comparison on is a separate question from whether the comparison is any good, and dev separates the two. Every action in the widened pool was judged, so each ordering can be replayed at matched budgets over one cache: judging the first $k$ actions in an order, applying the frozen rule, and scoring the result. At $k=2,6,12,18$ the gap-based order gives $-0.0009,-0.0016,-0.0054,-0.0057$, an AcuRank-style uncertainty order $-0.0013,-0.0048,-0.0057,-0.0057$, a random order $+0.0000,-0.0006,-0.0054,-0.0057$, and a cross-validated expected-gain model (ridge on $\log_2$ slot, gap, judge margin, margin/slot, order agreement; five query folds) $-0.0008,-0.0008,-0.0008,-0.0057$. Ordering the same judgements by their true effect, i.e.\ giving the allocation oracle, yields at most $+0.0003$: even a perfect budget allocation cannot rescue the stage. Perfect \emph{selection} among the judged actions would give $+0.0170$ at $k=2$ and $+0.0593$ at $k=18$. An independent control that judges a random subset of the pool reproduces the same picture ($-0.0010,-0.0018,-0.0035,-0.0109$ at the four budgets, 18{,}684 prompts) and is better than the gap order at $k=12$, so the allocation does not order the two policies consistently. The stage therefore loses the gain at the judgement, not at the allocation: the model prefers the wrong document roughly as often as the right one (AUC 0.532 for whether an action helps). This reframes where research effort belongs, and it is why the bounded CE-sourced pool, where the comparison is easy, is the only configuration that pays.

\paragraph{Condition-focused comparison.} If the judge cannot see the recoverable gain in the widened pool, one may ask whether it is being asked the wrong question. A second prompt keeps the candidate pool, the pairs, the two presentation orders and the decision rule, and changes only the question, from which passage is more relevant to which passage better satisfies the conditions the query states (time period or version, entity or jurisdiction, numeric bounds, exclusions). This arm is far more conservative: it applies 15 swaps instead of 51, of which 12 help and 3 hurt, and raises domain-macro nDCG@10 by $+0.0044$ (95\% CI $[+0.0012,+0.0085]$) at the same 9{,}342 prompts. Its precision (80\%) exceeds the relevance question's (63\%), but it finds a third as many swaps. The precision advantage holds on the other two starts: $+0.0071$ (CI $[+0.0024,+0.0123]$, 29 helpful and 10 harmful swaps) from the dense start and $+0.0034$ (CI $[-0.0003,+0.0074]$, 27 and 8) from the reranker start, so it is a property of the question rather than of the fused ranking. It is still smaller than the relevance question's gain on both starts ($+0.0106$ and $+0.0131$).  Asking for both in one prompt does not help: with the relevance question and the condition question merged into a single instruction the arm returns 0.3626 (\Delta +0.0045, CI [-0.0005, +0.0102]) with 31 helpful and 17 harmful swaps, i.e. it behaves like the relevance question and loses the precision that the narrow question buys. The two framings are therefore complementary rather than composable, and enforcing both requires two passes. The two questions are close to complementary, so a combined rule that applies the condition-based swap when the condition judge fires and the relevance-based swap otherwise reaches $0.3661$ ($+0.0080$, CI $[+0.0030,+0.0138]$, 54 swaps); the paired improvement over the relevance arm alone is $+0.0015$ (CI $[-0.0001,+0.0044]$), i.e.\ not established, and it costs twice the prompts. We report both arms and keep the single question as the deployed configuration.\\
The widened space is also not specific to the fused start: starting from the dense retriever alone, the same 18-action pool gives $+0.0026,+0.0012,+0.0011,+0.0006$ at the four budgets, against $+0.0106$ for the bounded tail protocol on that start. Widening costs most of the gain on both systems, which is the second-start evidence the acceptance criterion asks for.

\paragraph{Position of the correction.} The recoverable-gain analysis says that ranks 1--3 hold 43.6% of what a correct fix could return, so we re-ran the bounded protocol with the same pool, the same nine pairs per query and the same cost, but comparing each candidate with ranks 1--3 instead of 8--10. It loses: nDCG@10 0.3539 (\Delta -0.0042, CI [-0.0214, +0.0137]) with 23 helpful and 33 harmful swaps, against 0.3646 (\Delta +0.0065, CI [+0.0021, +0.0117]) with 32 helpful and 19 harmful for the tail window. The pool is not the problem -- the same documents are on offer -- so the difference has to come from the incumbent: in the tail the incumbent is often already doubtful and a swap is close to free, while overturning a confident top-ranked document needs evidence the pairwise question does not supply.

\paragraph{Two further attempts at a better judge.} Requiring an explicit evidence span before the verdict (the judge quotes the deciding sentence, then answers) collapses coverage: on a 39-query stratified sample it swaps 4 times (2 helpful, 2 harmful) for -0.0006, at roughly four times the cost of the single-token question, with 7% of outputs unparsable. Sending only the queries where the relevance arm is unsure (chosen margin below 0.9, 230 of 519) to the condition judge does not help either: it fires 13 times, all neutral, so the cascade reproduces the relevance arm exactly (0.3646) at 26.0 prompts per query. A second opinion is not independent where the first one is already unsure.

\paragraph{Transfer to FollowIR.} To check that the stage is not tied to the RETECO corpus we re-ran the whole frozen protocol, from dense encoding onward, on three FollowIR collections. Table~\ref{tab:followir} reports the graded nDCG@10 of the frozen system under the original and the changed instructions and the Direct judge's binary-gain delta. The frozen fusion beats both the dense and the cross-encoder-only ranking in all six instruction splits, so the starting system replicates; the replacement stage is positive on all three collections (+0.0084, +0.0072, +0.0064) though only Robust04's interval excludes zero. The effect is therefore not specific to RETECO, but the transfer evidence is weaker than the in-domain result. The official p-MRR tells the complementary story: under the changed instruction the frozen system scores $-0.0433$ on Core17, $-0.0208$ on News21 and $-0.0347$ on Robust04, and no system in any collection promotes a single newly relevant document. The positive nDCG deltas therefore measure ranking quality on the changed query text, not instruction following. Applying the judge's swaps under the changed instructions lifts graded nDCG@10 from 0.1761 to 0.1825 (Core17), 0.2849 to 0.2978 (News21) and 0.2382 to 0.2430 (Robust04), yet only 1, 5 and 2 of the 15, 21 and 28 applied swaps remove a document that the instruction change demoted: the stage improves the ranking without explicitly targeting the changed constraint, which is the gap the next phase has to close.

\begin{table}[t]
\centering
\scriptsize
\setlength{\tabcolsep}{3pt}
\begin{tabular}{lrrrrr}
\toprule
collection & topics & frozen (og) & frozen (chg) & Direct $\Delta$ & 95\% CI \\
\midrule
Core17 & 20 & 0.2674 & 0.1761 & +0.0084 & [-0.0017, +0.0186] \\
News21 & 32 & 0.3401 & 0.2849 & +0.0072 & [-0.0014, +0.0161] \\
Robust04 & 52 & 0.3092 & 0.2382 & +0.0064 & [+0.0018, +0.0114] \\
\bottomrule
\end{tabular}
\caption{FollowIR transfer. Graded nDCG@10 of the frozen system under the original and changed instructions; the last two columns are the Direct judge's delta in the harness binary-gain metric and its 95\% paired-bootstrap interval over system-query pairs.}
\label{tab:followir}
\end{table}

\textbf{Bidirectional agreement is load-bearing; abstention is not.} Judging each pair in one order only drops the gain to +0.0015 with an interval that includes zero, while removing the abstain option changes nothing (+0.0066). The benefit comes from requiring the judge to prefer the challenger in both presentation orders, not from letting it decline.

\textbf{Learning when to swap did not add anything.} We split the 1,211 official training queries by domain into 969 for the gate and 242 for calibration; three folds trained cross-encoders on two thirds of the 969 so that no ranking used for training a gate came from a model that had seen its query, and a fourth model trained on all 969 to rank the calibration queries. A logistic regression on rank features, dense and cross-encoder score gaps, rank disagreement, judge margin and order consistency, with its threshold chosen only on the 242 calibration queries, reproduces Direct exactly on dev (+0.0065 at tau = 0.0549). That threshold is not a filter, however: it coincides with the smallest eligible calibration probability, and the gate's calibration delta at that point is -0.0014, worse than the 0.0 of swapping nothing. Choosing the threshold on the calibration split alone (tau = 0.1026) gives +0.00015 on calibration and +0.00086 on dev. The learned gate thus neither removes errors nor saves calls, which is the negative result. Fixed margins of 0.7 and 0.9, a cross-encoder-gap rule, and a cheap pre-call gate that avoids 84\% of judge calls all score no better; a simple "cross-encoder or dense gap" ordering rule is in fact the strongest dev rule at +0.0072, which says the judge's value here lies in the comparison itself rather than in a learned policy over it. On the training side the picture is flatter than on dev: the three held-out folds give +0.0007 [-0.0033, +0.0041], +0.0032 [+0.0002, +0.0066] and +0.0068 [+0.0022, +0.0124], and the 242 calibration queries give -0.0014 [-0.0059, +0.0028]. Because those folds also differ in their cross-encoder and in their starting nDCG (0.2562, 0.2796 and 0.3030 against the frozen 0.3581), they show that the gain varies with the query split and the starting system; we report the measured dev gain and the other splits rather than converting them into a fresh-split expectation.

\textbf{Cost.} Direct issues 18 prompts per query at most (three candidates, three slots, two orders), 9,342 prompts and 5.31M prompt tokens over 519 queries, 873 seconds on one RTX 3090 with 11.2~GB peak, i.e.\ 1.7~s and 10.2k tokens per query - about 8.8 prompts per second. Context costs 1.77x tokens for no gain; per-document scoring costs less than Direct but also buys less.

\subsection*{9.1 Verification and replay}
\label{sec:verif}
Two audits changed what we report. First, the off-the-shelf reranker numbers in the arm table were computed with right-padded batches, so the logit read at the final position belonged to a pad token; a 26-query probe (two per domain) against an unpadded per-document reference gives a mean rank correlation of 0.43 and a maximum score difference of 9.7 before the fix, and 0.998 / 0.07 after it, and the corrected numbers are the ones reported here. The pairwise judge is unaffected: its batches are left-padded, and its predictions reproduce exactly. Second, every number in the replacement tables is now recomputed on a CPU from the committed artefacts (\texttt{scripts/replay\_replacement.py}): the script rebuilds each policy's ranking from \texttt{pair\_details} and the frozen top ten, re-derives the oracle and the four decision-rule ablations from the stored per-order probabilities, and exits non-zero if any value differs from the committed one by more than $10^{-6}$. It reproduces all macros, swap counts, point totals and bootstrap intervals in this section.

\section*{10 Conclusion}
A zero-shot dense stage plus a cross-encoder fine-tuned on the 1,071 questions with an in-pool positive gives +0.0332 domain-macro nDCG@10 over the same pipeline with the cross-encoder left at zero-shot (0.3249), on all 13 dev domains, stable across seeds (0.347-0.358). Two failures remain: an incomplete pool (1.9\% of questions have no gold in the dense top-1,000) and multi-evidence coverage (59.5\% of steps in the fused top ten). Every explicit temporal intervention - ten probes, one pre-registered test, step supervision, query reformulation, loss variants, backfilling, rank fusion - is neutral or negative. A bounded replacement stage over three external candidates is the one intervention that clears its pre-registered bar (0.3646, +0.0065, CI [+0.0021, +0.0117]), context around the compared documents adds nothing, and bidirectional judging is what makes the swap decisions reliable. Most transferable: a cross-encoder 0.04 better on its own metric can be worse fused, because fusion anchors on the first stage and re-normalises the second’s scores.

\begin{thebibliography}{99}
\small
\bibitem[Abdallah et al.(2026)]{abdallah2026tempo} Abdelrahman Abdallah, Mohammed Ali, Muhammad Abdul-Mageed, and Adam Jatowt. 2026. TEMPO: A Realistic Multi-Domain Benchmark for Temporal Reasoning-Intensive Retrieval. arXiv:2601.09523.
\bibitem[Abdallah et al.(2026)]{abdallah2026reteco} Abdelrahman Abdallah, Mohammed Ali, Muhammad Abdul-Mageed, Kevin Duh, and Adam Jatowt. 2026. RETECO: A SemEval-2027 Shared Task on Reasoning-Oriented Retrieval. Task proposal, University of Innsbruck. The organizers ask participants to cite TEMPO (and RECOR for Track 2) until the task description paper appears.
\bibitem[Robertson(2009)]{robertson2009bm25} Stephen Robertson and Hugo Zaragoza. 2009. The Probabilistic Relevance Framework: BM25 and Beyond. \emph{Foundations and Trends in Information Retrieval}, 3(4):333-389.
\bibitem[Nogueira(2019)]{nogueira2019passage} Rodrigo Nogueira and Kyunghyun Cho. 2019. Passage Re-ranking with BERT. arXiv:1901.04085.
\bibitem[Cormack et al.(2009)]{cormack2009rrf} Gordon V. Cormack, Charles L. A. Clarke, and Stefan Buettcher. 2009. Reciprocal Rank Fusion Outperforms Condorcet and Individual Rank Learning Methods. In \emph{SIGIR 2009}, pages 758-759.
\bibitem[MSR-TR-2010-82(2010)]{burges2010lambdarank} Christopher J. C. Burges. 2010. From RankNet to LambdaRank to LambdaMART: An Overview. Technical Report MSR-TR-2010-82, Microsoft Research.
\bibitem[Wang et al.(2018)]{wang2018lambdaloss} Xuanhui Wang, Cheng Li, Nadav Golbandi, Michael Bendersky, and Marc Najork. 2018. The LambdaLoss Framework for Ranking Metric Optimization. In \emph{CIKM 2018}, pages 1313-1322.
\bibitem[Jarvelin(2002)]{jarvelin2002ndcg} Kalervo Jarvelin and Jaana Kekalainen. 2002. Cumulated Gain-based Evaluation of IR Techniques. \emph{ACM Transactions on Information Systems}, 20(4):422-446.
\bibitem[Chen et al.(2024)]{chen2024m3} Jianlv Chen, Shitao Xiao, Peitian Zhang, Kun Luo, Defu Lian, and Zheng Liu. 2024. BGE M3-Embedding: Multi-Lingual, Multi-Functionality, Multi-Granularity Text Embeddings Through Self-Knowledge Distillation. arXiv:2402.03216.
\bibitem[card(2024)]{mxbai2024card} mixedbread.ai. 2024. mxbai-embed-large-v1 model card. https://huggingface.co/mixedbread-ai/mxbai-embed-large-v1 No peer-reviewed description exists for this checkpoint; we cite the model card and report the exact revision in \texttt{results/MODEL\_MANIFEST.json}.
\bibitem[Zhang et al.(2022)]{zhang2022hlatr} Yanzhao Zhang, Dingkun Long, Guangwei Xu, and Pengjun Xie. 2022. HLATR: Enhance Multi-Stage Text Retrieval with Hybrid List Aware Transformer Reranking. arXiv:2205.10569.
\bibitem[Min et al.(2021)]{mao2021jpr} Sewon Min, Kenton Lee, Ming-Wei Chang, Kristina Toutanova, and Hannaneh Hajishirzi. 2021. Joint Passage Ranking for Diverse Multi-Answer Retrieval. In \emph{EMNLP 2021}, pages 6997-7008.
\end{thebibliography}

\section*{Appendix A. The six-step ladder in full}
\subsection*{A.1 The loss form is not the bottleneck}
Three alternatives suggested by the coverage gap - averaging the log-loss over all positives, an nDCG-weighted permutation loss (LambdaLoss-style, cutoff 10), and domain-balanced (sqrt) sampling - were trained on the same 942 groups with the same hyper-parameters and scored on the same 45-query split:

\begin{table}[t]
\centering
\footnotesize
\resizebox{\columnwidth}{!}{%
\begin{tabular}{lrr}
\toprule
objective / sampler & best epoch & val domain-macro nDCG@10 \\
\midrule
current (log-sum-exp over positives) & 2 & 0.2865 \\
mean over positives & 2 & 0.2809 \\
nDCG-weighted permutations & 2 & 0.2764 \\
current + sqrt(domain) sampling & 2 & 0.2912 \\
zero-shot reference & - & 0.2730 \\
\bottomrule
\end{tabular}}%
\end{table}

Selecting on domain-macro rather than the plain mean picks the same epoch in all four arms, so the official aggregation does not move the checkpoint; the split cannot resolve the gaps (+-0.03, S7), so only sqrt sampling’s +0.0047 ([-0.0164, +0.0285]) is even measurable, and the pre-registered 0.010 tie band keeps the current objective. Sqrt sampling’s one visible effect is raising the share of queries with a relevant document in the top ten from 93.3\% to 95.6\% without moving nDCG: it makes the easy decision easier.

Dev measures the incumbent itself. Retrained through the current code path (same 942 groups, two epochs, its own val-preferred epoch) it scores \textbf{0.3575}, 0.0006 below the frozen artifact (domain-clustered CI [-0.0037, +0.0021]), so re-running the recipe is free within measurement error and the frozen artifact is a draw from the same band. The three replacement loss forms stay unmeasured on dev - the screen separates them by at most 0.015, below the dev half-width of 0.007-0.014 - so dev cannot say they lose; what it can say is that the screen ranks arms at 0.2764-0.2912, cannot resolve 0.003, and that the arm it did put on dev reproduces the frozen number.

\subsection*{A.2 Forcing every evidence step into the top ten does not pay}
If the fused top ten covers only 59.5\% of annotated steps (51.4\% for the cross-encoder alone), the mechanical remedy is to require every step to have a positive ranked high. We trained the step-augmented objective of S3 at two weights and scored the systems on dev with the frozen fusion:

\begin{table*}[t]
\centering
\scriptsize
\resizebox{\textwidth}{!}{%
\begin{tabular}{lrrrrr}
\toprule
system & dev domain-macro & Delta vs frozen & domain-clustered 95\% CI & steps in top 10 & questions with 2+ relevant in top 10 \\
\midrule
frozen (no step term) & 0.3581 & - & - & 59.45\% & 38.31\% \\
L\_step, lambda = 0.3 & 0.3593 & +0.0012 & [-0.0021, +0.0048] & 59.54\% & 38.20\% \\
L\_step, lambda = 1.0 & 0.3600 & +0.0019 & [-0.0014, +0.0058] & 59.69\% & 38.47\% \\
\bottomrule
\end{tabular}}%
\end{table*}

The coverage column counts the fused list, whose top ten does not change when the list is truncated from 1000 to 100 rows (59.45\%); the cross-encoder’s own ordering over the same candidates covers 51.39\% before the term and 48.14\% / 45.78\% with it, so the degradation is a property of the ordering, not of depth.

The term does not pay and does not move the quantity it optimises: coverage stays at 59.5\% / 59.7\%. One reading fits: a step’s positives are mostly documents the ordinary listwise term already covers (3.44 relevant documents over 2.34 steps, so a step rarely adds a \emph{different} document), and the log-sum-exp inside a step is dominated by the step whose positive is hardest to retrieve, so the gradient chases documents that are absent or buried. The annotations are not at fault: every one of the 2,873 step-level positives is also a query-level positive. The instructive part is the selection: on the internal split lambda = 1.0 was the \emph{best} of six arms trained that day (0.2913, above 0.2865), and on dev the two step arms are indistinguishable from the frozen system - which is why S3 fixes the fusion form in advance and S7 refuses to select on 45 questions.

\subsection*{A.3 Two further corrections that also lose}
The frozen system feeds the reranker a title-only query and excludes the 140 training queries whose dense top-100 contains no positive. Both were corrected with the frozen recipe and scored on dev with the frozen fusion:

\begin{table}[t]
\centering
\footnotesize
\resizebox{\columnwidth}{!}{%
\begin{tabular}{lrrrr}
\toprule
system & dev domain-macro & Delta vs frozen & query-clustered 95\% CI & domain-clustered 95\% CI \\
\midrule
frozen & 0.3581 & - & - & - \\
full-question query, 64-word cap & 0.3443 & -0.0138 & [-0.0272, -0.0005] & [-0.0338, +0.0052] \\
+ 140 backfilled queries, 294 positives & 0.3565 & -0.0017 & [-0.0088, +0.0061] & [-0.0061, +0.0061] \\
\bottomrule
\end{tabular}}%
\end{table}

The two no longer cost the same: the backfill is a wash and the query rewrite pays. For the query arm the loss is smaller where the body conditions live (S5.3), so the extra text carries signal but not enough to pay for the harm on the other 278 queries. For the backfill arm the 294 restored positives are, by construction, documents the first stage never retrieves, so training on them spends capacity on a decision never made at inference, where the candidate list is the dense top-100 - and on dev that costs nothing measurable. The unchanged objective re-run through the current code path scores 0.3575, 0.0006 below the frozen artifact, so the full-question arm loses 0.0132 against its own baseline and the step arms sit 0.0018 / 0.0025 above it. Dev therefore supports a deployment decision - six of the ladder’s steps retrained, none showing a reliable gain over it, and the sixth (dense pooling) a wash at +0.0014 - plus a negative result against a same-code baseline that the 45-query screen cannot rank at all.

\subsection*{A.4 Where the gain sits: not on the multi-step questions}
Partitioning dev by the number of official evidence steps a question needs (\texttt{results/STEP\_SPLIT.md}), the frozen system’s gain over zero-shot CE is largest on single-step questions and weakest on two-step ones:

\begin{table*}[t]
\centering
\scriptsize
\resizebox{\textwidth}{!}{%
\begin{tabular}{lrrrrr}
\toprule
partition & queries & zero-shot CE & frozen SFT & delta & 95\% CI (domain-clustered) \\
\midrule
1 step & 75 & 0.3185 & 0.3896 & +0.0711 & [+0.0244, +0.1056] \\
2 steps & 151 & 0.3857 & 0.4058 & +0.0200 & [-0.0033, +0.0463] \\
3+ steps & 258 & 0.3196 & 0.3511 & +0.0315 & [+0.0084, +0.0571] \\
none & 35 & 0.1579 & 0.1705 & +0.0126 & [+0.0030, +0.0524] \\
\bottomrule
\end{tabular}}%
\end{table*}

The objective buys a better \emph{single-document} decision almost everywhere but does not teach evidence assembly: on the 151 two-step questions the gain is within noise. That is the gap the step term targets, and it fixes the bar those arms have to clear.

\end{document}
